% ============================================================
%  Příprava na maturitu – Mechatronika
%  Kompiluj: pdflatex maturita_mechatronika.tex
%  Kódování: UTF-8
% ============================================================

\documentclass[a4paper,10pt]{article}

% --- Balíčky ---
\usepackage[T1]{fontenc}
\usepackage[utf8]{inputenc}
\usepackage[czech]{babel}
\usepackage[left=2.4cm, right=1.8cm, top=1.8cm, bottom=1.8cm]{geometry}
\usepackage{lmodern}
\usepackage{microtype}
\usepackage{parskip}
\usepackage{titlesec}
\usepackage{enumitem}
\usepackage{xcolor}
\usepackage{hyperref}
\usepackage{tabularx}
\usepackage{booktabs}
\usepackage{array}
\usepackage{tikz}
\usepackage{eso-pic}
\usepackage{mdframed}
\usepackage{amsmath}
\usepackage{amssymb}
\usepackage{pgfplots}
\pgfplotsset{compat=1.18}
\usetikzlibrary{arrows.meta, positioning, shapes.geometric, calc}

% --- Barvy ---
\definecolor{accent}{HTML}{03254E}
\definecolor{stripeblue}{HTML}{568EA3}
\definecolor{medgray}{HTML}{888888}
\definecolor{lightblue}{HTML}{EAF2F8}
\definecolor{lightyellow}{HTML}{FDFBE4}
\definecolor{lightgreen}{HTML}{EAF7EA}

% --- Modrý pruh vlevo ---
\AddToShipoutPictureBG{%
  \begin{tikzpicture}[remember picture, overlay]
    \fill[stripeblue]
      (current page.north west)
      rectangle
      ([xshift=10mm] current page.south west);
    \fill[accent!60]
      ([xshift=10mm] current page.north west)
      rectangle
      ([xshift=12mm] current page.south west);
  \end{tikzpicture}%
}

\hypersetup{colorlinks=true, urlcolor=accent, linkcolor=accent}

\titleformat{\section}
  {\large\bfseries\color{accent}}
  {\thesection.}
  {0.5em}
  {}
  [\color{accent}\titlerule]

\titlespacing{\section}{0pt}{16pt}{6pt}

\newmdenv[
  backgroundcolor=lightblue,
  linecolor=stripeblue,
  linewidth=1pt,
  roundcorner=4pt,
  innerleftmargin=8pt,
  innerrightmargin=8pt,
  innertopmargin=6pt,
  innerbottommargin=6pt,
  skipabove=4pt,
  skipbelow=4pt
]{pojmybox}

\newmdenv[
  backgroundcolor=lightyellow,
  linecolor=accent!40,
  linewidth=1pt,
  roundcorner=4pt,
  innerleftmargin=8pt,
  innerrightmargin=8pt,
  innertopmargin=6pt,
  innerbottommargin=6pt,
  skipabove=4pt,
  skipbelow=8pt
]{vysvetlenibox}

\newmdenv[
  backgroundcolor=lightgreen,
  linecolor=accent!30,
  linewidth=1pt,
  roundcorner=4pt,
  innerleftmargin=8pt,
  innerrightmargin=8pt,
  innertopmargin=6pt,
  innerbottommargin=6pt,
  skipabove=4pt,
  skipbelow=8pt
]{vzorcebox}

\pagestyle{plain}

% ============================================================
\begin{document}

% --- Záhlaví ---
\begin{minipage}[t]{0.75\textwidth}
    {\Huge\bfseries\color{accent} Příprava na maturitu}\\[4pt]
    {\large\color{stripeblue} Mechatronika}\\[2pt]
    \textcolor{medgray}{\small Suchánek Lukáš \quad SPŠ Prosek \quad 2026}
\end{minipage}
\hfill
\begin{minipage}[t]{0.22\textwidth}
    \raggedleft
    \begin{tikzpicture}
        \draw[color=stripeblue, rounded corners=4pt, line width=1pt]
            (0,0) rectangle (3,1.2)
            node[midway, align=center, color=accent, font=\small\bfseries]
            {Otázek: 22};
    \end{tikzpicture}
\end{minipage}

\vspace{6pt}
\noindent\color{accent}\rule{\linewidth}{1.5pt}
\vspace{2pt}

% ============================================================
\section{Umělá inteligence}

\begin{pojmybox}
\textbf{Klíčové pojmy:}
strojové učení, neuronová síť, hluboké učení, CNN, RNN, LSTM,
regrese, klasifikace, trénovací/validační/testovací data,
overfitting, přeučení, aktivační funkce, zpětná propagace
\end{pojmybox}

\begin{vysvetlenibox}
\textbf{Rozdělení metod AI}

\begin{center}
\begin{tikzpicture}[scale=0.85,
  node/.style={rectangle, draw=accent, fill=lightblue, rounded corners=3pt,
               minimum width=3cm, minimum height=0.55cm,
               text centered, font=\small}]
  \node[node, fill=lightyellow] (ai) at (0,3) {Umělá inteligence};
  \node[node] (ml) at (-3,2)  {Strojové učení};
  \node[node] (es) at (0,2)   {Expert. systémy};
  \node[node] (cv) at (3,2)   {Poč. vidění};
  \node[node] (sup)  at (-4.5,1)  {Superv.};
  \node[node] (uns)  at (-3,1)    {Nesuperv.};
  \node[node] (rl)   at (-1.5,1)  {Posilovací};
  \node[node] (dl)   at (-3,0)    {Hluboké učení};
  \node[node] (cnn)  at (-4.2,-0.8) {CNN};
  \node[node] (rnn)  at (-3,-0.8)   {RNN/LSTM};
  \node[node] (tr)   at (-1.8,-0.8) {Transformer};
  \draw[>-,accent] (ai)--(ml);
  \draw[>-,accent] (ai)--(es);
  \draw[>-,accent] (ai)--(cv);
  \draw[>-,accent] (ml)--(sup);
  \draw[>-,accent] (ml)--(uns);
  \draw[>-,accent] (ml)--(rl);
  \draw[>-,accent] (sup)--(dl);
  \draw[>-,accent] (dl)--(cnn);
  \draw[>-,accent] (dl)--(rnn);
  \draw[>-,accent] (dl)--(tr);
\end{tikzpicture}
\end{center}

\textbf{Neuronová síť -- princip}

\begin{vzorcebox}
Jeden neuron: $y = f\!\left(\sum_{i=1}^{n} w_i x_i + b\right)$\\
Aktivační funkce: ReLU: $f(x) = \max(0,x)$; \quad
Sigmoid: $f(x) = \frac{1}{1+e^{-x}}$\\
Zpětná propagace: $w \leftarrow w - \eta \frac{\partial L}{\partial w}$
\quad ($\eta$ = rychlost učení, $L$ = ztrátová funkce)
\end{vzorcebox}

\textbf{Použití AI v mechatronice}

\begin{center}
\begin{tabularx}{0.95\linewidth}{@{} l X l @{}}
    \toprule
    \textbf{Aplikace} & \textbf{Popis} & \textbf{Metoda} \\
    \midrule
    Detekce vad    & Rozpoznání defektů na výrobku z kamery & CNN \\
    Predikce poruch & Předpověď poruchy z vibrací a teplot & LSTM \\
    Autonomní robot & Navigace, vyhýbání překážkám & Posilovací učení \\
    Kvalita svaru  & Klasifikace kvality z obrazu & CNN \\
    Digitální dvojče & Aktualizace modelu dle měření & Hybridní \\
    \bottomrule
\end{tabularx}
\end{center}
\end{vysvetlenibox}

% ============================================================

% ============================================================
\end{document}
